\section{Related Work}
\label{sec:related}

\para{Human route language and vision-language navigation.}
\rtr introduced a large benchmark of human-written indoor navigation instructions grounded in real Matterport3D environments \cite{anderson2018vision}. This makes it a useful source of natural human route language: workers describe room transitions, turns, landmarks, and stopping points in task-directed prose. Speaker-Follower models later used generated instructions to improve \vln agents, showing that synthetic route language can affect downstream navigation behavior \cite{fried2018speaker}. \textsc{Room-Across-Room} extends this line with multilingual, densely grounded route descriptions \cite{ku2020room}. Building on this work, we use \rtr not to train a navigation agent, but to ask whether model-written spatial stories preserve a different linguistic profile from human route descriptions.

\para{\llm spatial reasoning and spatial representation.}
Recent work suggests that spatial language fluency does not guarantee robust spatial reasoning. Visualization-of-Thought prompting improves \llm performance on navigation and grid-based spatial tasks, but the need for explicit state visualization highlights the difficulty of spatial reasoning from text alone \cite{wu2024mindseye}. Tag Map shows a complementary approach: make scene context explicit as a text-based map that an \llm can query for grounded planning \cite{zhang2024tagmap}. Other work compares \llm spatial preposition behavior against human spatial-representation data and finds both human-like vertical preferences and non-human axis-specific patterns \cite{wu2025spatial}. Hybrid Mind similarly frames spatial cognition in terms of landmark, route, and survey knowledge, finding weak unaided \llm performance and stronger tool-augmented performance \cite{yang2025spatialcognition}. Our study differs by measuring the textual structure of spatial stories, not task-answer correctness.

\para{Human-vs-\llm style and story generation.}
Human-vs-\llm writing studies show that generated text can be separated from human text using grammatical and rhetorical features. Reinhart et al. find systematic differences across Biber-style features and show that instruction-tuned models diverge more from human style than base models \cite{reinhart2025llms}. Story generation work provides human prose resources and evaluation concerns relevant to our control corpus. \writingprompts pairs prompts with large-scale human-written stories \cite{fan2018hierarchical}, while StoryAlign argues that story evaluation should track human preferences over coherence, creativity, characterization, fluency, and relevance \cite{xia2026storyalign}. A recent narrative-theory survey also argues against a single generic narrative-quality benchmark and favors theory-specific attributes \cite{liu2026narrative}. We follow that recommendation by focusing on a narrow attribute family: spatial movement through built environments.
